% 20_GUIA_REPORTE.tex — guía de lectura del informe H7 (ADR-59 y adendas 1 y 2)
% Compilar desde docs/:
%   python ../scripts/12_reporte_html.py --out ../reporte.html --tex cifras_reporte.tex
%   pdflatex 20_GUIA_REPORTE.tex && pdflatex 20_GUIA_REPORTE.tex
% Ninguna cifra de resultados se teclea: vienen de cifras_reporte.tex (\cifra...).
\documentclass[11pt,a4paper]{article}
\usepackage[T1]{fontenc}
\usepackage[utf8]{inputenc}
\IfFileExists{spanish.ldf}{\usepackage[spanish,es-noshorthands,es-nodecimaldot]{babel}}{}
\usepackage[scaled=0.95]{helvet}
\renewcommand{\familydefault}{\sfdefault}
\usepackage{amsmath,amssymb}
\usepackage[margin=2.4cm]{geometry}
\usepackage{booktabs,tabularx,enumitem,microtype}
\usepackage[hidelinks]{hyperref}
\usepackage{fancyhdr}
\setlength{\parindent}{0pt}
\setlength{\parskip}{0.55em}
\setlist{itemsep=0.2em,topsep=0.3em}
\renewcommand{\contentsname}{Índice}
\renewcommand{\tablename}{Tabla}
\pagestyle{fancy}\fancyhf{}
\fancyhead[L]{\small Guía de lectura · La historia de un entrenador}
\fancyhead[R]{\small\thepage}
\renewcommand{\headrulewidth}{0.4pt}

\InputIfFileExists{cifras_reporte.tex}{}{%
  \typeout{FALTA cifras_reporte.tex: genera las macros con 12_reporte_html.py --tex}%
  \newcommand{\cifraPosJ}{\textbf{??}}\newcommand{\cifraPosJic}{\textbf{??}}%
  \newcommand{\cifraSolJ}{\textbf{??}}\newcommand{\cifraSolJic}{\textbf{??}}%
  \newcommand{\cifraSolJq}{\textbf{??}}\newcommand{\cifraOrtJ}{\textbf{??}}%
  \newcommand{\cifraOrtJic}{\textbf{??}}\newcommand{\cifraOrtJq}{\textbf{??}}%
  \newcommand{\cifraOrtJmargen}{\textbf{??}}\newcommand{\cifraFtJ}{\textbf{??}}%
  \newcommand{\cifraFtJic}{\textbf{??}}\newcommand{\cifraPresJO}{\textbf{??}}%
  \newcommand{\cifraPresJOic}{\textbf{??}}\newcommand{\cifraPresMargen}{\textbf{??}}%
  \newcommand{\cifraCtxRech}{\textbf{??}}\newcommand{\cifraRotJ}{\textbf{??}}%
  \newcommand{\cifraDauc}{\textbf{??}}\newcommand{\cifraViajan}{\textbf{??}}%
  \newcommand{\cifraMarcador}{\textbf{??}}\newcommand{\cifraFirmaCru}{\textbf{??}}%
  \newcommand{\cifraFirmaDid}{\textbf{??}}}

\newcommand{\nivel}[1]{\fbox{\footnotesize\textbf{#1}}}
\newcommand{\fuente}[1]{{\small\ttfamily #1}}
\newcommand{\bloque}[3]{\par\textbf{Qué muestra.} #1\par\textbf{Cómo se lee.} #2\par\textbf{Sustento.} #3\par}

\title{\textbf{La historia de un entrenador}\\[0.3em]\large Guía de lectura del informe interactivo}
\author{dt-decoder · Hackathon 2026}
\date{Acompaña a \texttt{reporte.html} (ADR-59 y sus adendas 1 y 2)}

\begin{document}
\maketitle
\thispagestyle{empty}

\begin{abstract}
\noindent Esta guía acompaña a la página \texttt{reporte.html}. La página cuenta la
historia; la guía dice qué significa cada parte, cómo interpretarla y en qué
matemática descansa. Las dos se generan desde los mismos siete archivos JSON:
toda cifra de resultados de esta guía sale de ellos, igual que en la página. Los
números que quedan escritos a mano son definiciones y umbrales fijados antes de medir.
\end{abstract}

\tableofcontents
\newpage

%======================================================================
\section{Cómo usar la página}

La página cuenta cinco historias (Jardine, Larcamón, Ambriz, Herrera y Ortiz)
en tres actos (adenda 2 de ADR-59). El \textbf{acto 1} es común: cómo se lee
una posesión y en qué confiar. El \textbf{acto 2} describe la era principal del
técnico (la de más partidos) contra la liga del mismo torneo. El \textbf{acto 3}
mira el club antes y después de él, y qué cambia cuando cambia de club. El
\textbf{cierre} pone a prueba el método.

La barra superior tiene tres controles: el selector de historia, el índice por
actos y el interruptor \emph{sencilla / técnica}. En modo sencillo se ocultan
los intervalos, el valor de $q$ y los plegables de método; la lectura de $q$ en
palabras, el nivel de cada frase y el margen de cada nulo se ven siempre. La
portada de cada historia trae cinco frases en ranuras fijas, cada una con un
enlace a su sección.

Cada sección de los actos 2 y 3 tiene cuatro capas: la frase con su color, la
figura, una traducción ``en un partido'' y el plegable ``cómo lo medimos''. Si
falta una capa, la página lo dice. En modo técnico, arriba de todo están las
huellas de los insumos: el commit del código y, para cada JSON, la fecha de su
corrida y el inicio de su sha256.

Esta guía detalla la historia de Jardine, que abre la página. Las otras cuatro
usan las mismas plantillas de redacción (sección~\ref{sec:adenda2}).

Tres gestos bastan para leerla:
\begin{itemize}
\item \textbf{Tocar una cifra subrayada} muestra el archivo y el campo de donde sale.
\item \textbf{Los selectores de cada figura} (corregido/crudo, métrica, técnico)
  redibujan solo esa figura.
\item \textbf{Pasar el cursor o tocar} cualquier punto, barra o casilla muestra su valor exacto.
\end{itemize}

Si falta un JSON, la sección correspondiente no inventa nada: muestra el
comando que lo genera. Si los datos cambian y una frase preinscrita deja de
sostenerse, el generador se detiene con un mensaje y no escribe la página.

\subsection{Los tres niveles de evidencia}
Cada frase lleva una etiqueta y una sola. En la página, la etiqueta tiene
color: verde para A, ámbar para B, gris para C.

\begin{tabularx}{\linewidth}{@{}lXX@{}}
\toprule
nivel & qué es & cómo se redacta\\
\midrule
\nivel{A} probado & contraste escrito antes de medir, IC por bootstrap, sigue en pie tras la corrección de Benjamini--Hochberg dentro de su familia & ``difiere'', con IC y $q$\\
\nivel{B} medido & diferencia contra la liga del mismo torneo, con IC, fuera de una familia & ``por encima / por debajo de la liga'', con IC\\
\nivel{C} descriptivo & percentiles, perfiles, conteos, comparaciones entre clubes & ``en $X$ de $Y$ torneos\ldots'', sin IC\\
\bottomrule
\end{tabularx}

Un \textbf{nulo} lleva la etiqueta A con la marca ``nulo''. No dice que dos cosas
sean iguales: dice el tamaño máximo de diferencia que los datos permiten
descartar (sección~\ref{sec:nulos}).

La etiqueta \textbf{adenda 1} señala una frase cuyo nivel o regla cambió respecto al
borrador de ADR-59, siempre por la misma razón: el JSON no trae intervalo, o la
regla no estaba fijada. El detalle está en la sección~\ref{sec:adenda}.

%======================================================================
\section{Fundamentos comunes}

\subsection{Unidad de análisis}
\begin{itemize}
\item \textbf{Posesión}: la cadena de acciones de un equipo desde que recupera el
  balón hasta que lo pierde, sale o remata.
\item \textbf{Acción}: pase, conducción o remate.
\item \textbf{Zona}: la cancha partida en una malla de $5\times 4$, fijada antes de
  medir. El equipo analizado ataca siempre hacia la derecha.
\item \textbf{Fase}: juego abierto, transición, reanudación o balón parado.
  \textbf{Estado} = zona $\times$ fase.
\item \textbf{Era}: el par (club, entrenador). La clave es compuesta: Jardine en el
  América y Jardine en San Luis son dos eras.
\end{itemize}

\subsection{La cadena de Markov absorbente}
Cada acción mueve la posesión de un estado transitorio $i$ a otro $j$ con
probabilidad $Q_{ij}$, o la termina en un estado absorbente $k$ (remate, gol,
pérdida) con probabilidad $R_{ik}$. Con la matriz fundamental
\[
N = (I - Q)^{-1},
\]
el número esperado de acciones antes de terminar es $E[T] = N\mathbf{1}$ y la
probabilidad de terminar en cada absorbente es $B = N R$. Así se responden las
tres preguntas de la página: dónde (los estados que visita), cuánto dura
($E[T]$) y cómo termina ($B$).

\textbf{Límite declarado.} La distribución observada de longitudes no es la
geométrica que implica una cadena de primer orden; la cadena se usa para
esperanzas y probabilidades de absorción, no para simular (ADR-21).

\subsection{Deriva del proveedor y comparación contra la liga}\label{sec:deriva}
El proveedor cambió su forma de registrar eventos entre torneos, y el número
de acciones por posesión se movió para todos los clubes a la vez. Una resta
directa entre dos eras de un mismo club mezcla estilo y deriva. Sin corregir,
\cifraFirmaCru{} pares que difieren dan al técnico más reciente las posesiones
más largas; con la corrección, \cifraFirmaDid{}.

Por eso cada era $u$ se compara con una \textbf{base}: la liga sin los partidos
del club, en los torneos de la era, ponderada a su mezcla de torneos
($w_t$ = fracción de partidos de la era en el torneo $t$):
\[
D_u = \log E[T]\big(\hat P_u\big) - \log E[T]\big(\hat P_{\text{base}(u)}\big),
\qquad \text{rel}_u = e^{D_u} - 1 .
\]
Para dos eras $a,b$ del mismo club, $\theta_{ab} = D_a - D_b$, que se reporta
como $e^{\theta}-1$. No es una diferencia en diferencias clásica: cada era se
observa solo en su periodo. El supuesto es que, sin cambio de entrenador, el
club se habría movido como la liga.

Las métricas del proveedor (xG, OBV, pases progresivos, field tilt) siguen la
misma lógica en escala aditiva: diferencia = media de la era $-$ base, con
$\text{base} = \sum_t w_t\,\bar x_{\text{liga},t}$.

\subsection{Incertidumbre: bootstrap por partido}
Los partidos, no las acciones, son la unidad que se remuestrea: así se conserva
la dependencia entre acciones del mismo partido. En cada réplica se remuestrean
los partidos de la era y, estratificados por torneo, los de la base (compartida
dentro del club). El intervalo sale de los cuantiles de las réplicas; el tipo
(básico o percentil) lo declara cada JSON. El intervalo básico al 95\% es
\[
\big[\,2\hat\theta - \theta^*_{0.975},\; 2\hat\theta - \theta^*_{0.025}\,\big].
\]

\subsection{Comparaciones múltiples: el valor $q$}
Cuando se hacen $m$ contrastes, algunos saldrán ``distintos'' por azar. La
corrección de Benjamini--Hochberg controla la proporción esperada de falsos
descubrimientos. Con los valores $p$ ordenados $p_{(1)}\le\dots\le p_{(m)}$,
\[
q_{(i)} = \min_{j \ge i} \frac{m\,p_{(j)}}{j},
\]
y un contraste \textbf{sobrevive} si $q \le 0.05$. La familia (qué contrastes
entran en $m$) se declara antes de medir. La página muestra $q$ y nunca $p$,
acompañado de una lectura en palabras: ``muy sólido'' ($q\le 0.01$), ``sólido,
con margen estrecho'' ($q\le 0.05$), ``insuficiente tras corregir''
($q\le 0.15$), ``no lo distinguimos del azar''.

\subsection{Cómo se escribe un nulo}\label{sec:nulos}
Si el intervalo es $[L,U]$ y la diferencia no sobrevive, la página escribe
``no detectamos una diferencia mayor a $X$'', con
\[
X = \max\big(|L|,\,|U|\big),
\]
el extremo más lejano al cero. Si el intervalo no toca el cero pero la
diferencia no sobrevive a la corrección, la frase lo dice así, y $X$ se lee
como cota: ``si existe, es menor a $X$''.

%======================================================================
\section{Sección por sección}

\subsection*{Acto 1 · El marco \hfill{\normalfont\small componente 5.6}}
\bloque{El esquema de las tres preguntas y la serie de acciones por posesión de
la liga, torneo a torneo, con una banda de $\pm$ una desviación estándar entre
clubes.}
{La línea sube para toda la liga a la vez: eso es la deriva, no el estilo de un
equipo. Es la razón de comparar contra la liga del mismo torneo.}
{Sección~\ref{sec:deriva}. Fuente: \fuente{deriva\_proveedor › por\_torneo} y
\fuente{did\_h4\_v1 › firma\_temporal}. El bloque plegable resume la cadena.}

\subsection*{2.1, 2.3 y 3.1 · Con el balón y el club antes y después \hfill{\normalfont\small 5.1 ofensiva}}
\bloque{Duración de la posesión contra la liga (\nivel{B} \cifraPosJ{}
\cifraPosJic), contra las otras eras del club (\nivel{A} Solari \cifraSolJ{}
\cifraSolJic, $q=$ \cifraSolJq; nulo contra Ortiz \cifraOrtJ{} \cifraOrtJic,
$q=$ \cifraOrtJq), progresión, ocasiones, OBV y field tilt
(\nivel{B} \cifraFtJ{} \cifraFtJic). La figura de pares enfrenta el valor
corregido (punto lleno) con la resta cruda (aro).}
{Un punto a la derecha de la línea punteada es una posesión más larga que la
de la otra era. Si el intervalo corregido no toca la línea y $q\le 0.05$, la
frase es A. Contra Ortiz el intervalo no toca el cero, pero $q>0.05$: se
redacta como nulo con cota \cifraOrtJmargen.}
{\begin{itemize}
\item Duración: $\text{rel}_u$ y $\theta_{ab}$ de la sección~\ref{sec:deriva} (ADR-53).
\item npxG: suma por partido de \texttt{shot\_statsbomb\_xg} sin penales.
\item OBV: suma de \texttt{obv\_total\_net} del equipo.
\item Pase progresivo: completado, de juego, con $x_0 \ge 48$ y
  $d(x_1,y_1) \le 0.75\, d(x_0,y_0)$, donde $d(x,y)=\sqrt{(120-x)^2+(40-y)^2}$.
\item Field tilt: acciones del equipo con $x\ge 80$ en su marco, divididas entre
  las de ambos equipos con $x\ge 80$ en su propio marco. La media de la liga es
  $0.5$ por construcción.
\item El mapa de zonas reparte las acciones de la cadena en las 20 casillas; los
  porcentajes se redondean por restos mayores para que sumen 100.
\end{itemize}
Fuentes: \fuente{did\_h4\_v1}, \fuente{metricas\_v1 › global}.}

\subsection*{2.4 · Sin el balón \hfill{\normalfont\small 5.1 defensiva}}
\bloque{npxG concedido contra la liga (\nivel{B}) y el nivel de presión entre las
tres eras del América (\nivel{A}), con E4 como sustituto declarado de las
transiciones.}
{Contra Ortiz, \cifraPresJO{} \cifraPresJOic: no detectamos una diferencia de
nivel de presión mayor a \cifraPresMargen. El contraste Ortiz--Solari, que sí
sobrevive, muestra que el instrumento puede detectar diferencias dentro del
mismo club.}
{E2 = nivel de presión por acción del rival en la vista defensora, comparado
con la liga sin el club, estandarizada por torneo $\times$ zona (ADR-54).
E4 = posesiones rivales de una sola acción en juego abierto. Toda cifra de
presión es \emph{por acción}, no por posesión: los dos estimandos pueden
diferir en signo. Fuente: \fuente{did\_presion\_v1 › pares}.}

\subsection*{2.5 · ¿La misma idea torneo tras torneo? \hfill{\normalfont\small 5.2 consistencia}}
\bloque{La posesión relativa torneo a torneo y los percentiles de field tilt y
npxG a favor. El selector cambia la métrica.}
{En la vista de posesión, la línea punteada es la liga del torneo. En las
vistas de percentil, $0.5$ es la mediana de la liga; $1$ es el valor más alto.
Un torneo parcial (menos de 12 partidos de la era) no lleva percentil.}
{Percentil: la media de la era en el torneo dentro de la distribución de las
medias de los demás equipos. Nivel \nivel{C}: la serie por torneo no trae
intervalo (adenda 1). Fuentes: \fuente{did\_h4\_v1 › serie\_por\_torneo},
\fuente{metricas\_v1 › por\_torneo}.}

\subsection*{2.6 · ¿Ajusta al contexto? \hfill{\normalfont\small 5.2 variabilidad}}
\bloque{Lo que ajusta la liga entera en cuatro contextos (localía, marcador,
minuto 60, rival) y el bosque de $\theta$ de cada era del América.}
{Cada punto es cuánto más (o menos) ajusta la era que la liga. Aro: el
intervalo excluye el cero antes de corregir, pero no sobrevive. En el América,
\cifraCtxRech{} contrastes sobreviven.}
{Para un contexto con estados $A$ y $B$ y una métrica $M$,
\[
\theta = \big[M_{\text{era}}(A) - M_{\text{era}}(B)\big] - \big[M_{\text{liga}}(A) - M_{\text{liga}}(B)\big],
\]
desde la perspectiva del club (ADR-56). M1 = acciones por posesión, M2 =
P(remate) a favor, M3 = presión del que defiende, M4 = P(remate) concedido.
La familia de casos (ADR-57) se corrige por separado. La tabla de la liga es
\nivel{C} porque el JSON no trae su intervalo. Fuente: \fuente{contexto\_v1}.}

\subsection*{2.8 · Jugadores y minutos \hfill{\normalfont\small 5.3}}
\bloque{Rotación (continuidad del once y N80), los cinco jugadores con más
minutos y lo que pasa tras el primer cambio.}
{Cada anillo es un torneo y marca el percentil de la era en la liga; la marca
en el aro es el umbral (15\% inferior para la continuidad, mediana para N80).
La continuidad quedó en el 15\% inferior en \cifraRotJ{} torneos. Tocar un
jugador dibuja su mapa de zonas.}
{\begin{itemize}
\item N80: el menor número de jugadores que suman el 80\% de los minutos.
\item Continuidad: fracción media del once que se repite respecto al partido
  anterior del mismo torneo.
\item Tras el primer cambio táctico entre los minutos 55 y 80: ventanas de diez
  minutos antes y después, con una exclusión de $\pm$60 s; $\Delta$ = después
  $-$ antes; $\theta$ = media de $\Delta$ de la era $-$ media de la liga,
  ponderada por bloque de minuto, marcador y torneo (ADR-58).
\item La redacción es siempre ``tras sus primeros cambios, el equipo\ldots'': el
  técnico cambia cuando el partido lo pide, así que no hay lectura causal.
\end{itemize}
Fuente: \fuente{jugadores\_v1}. Los jugadores aparecen por identificador; el
JSON no trae nombres.}

\subsection*{2.7 · Balón parado \hfill{\normalfont\small 5.4}}
\bloque{El embudo de la liga (córner, tiro libre indirecto, saque de banda), lo
que pesa en el xG de un remate, el contraste del América en córners y las zonas
de remate a favor y en contra.}
{El embudo es la probabilidad de que la secuencia termine en remate o en gol.
En el bosque de coeficientes, a la izquierda del cero hay menos xG.}
{C1 $=P(\text{remate}\mid\text{secuencia})$. El modelo de xG es logístico:
\[
\operatorname{logit} P(\text{gol}) = \beta_0 + \beta_1\,\text{distancia} + \beta_2\,\text{ángulo} + \beta_3\,\text{cabeza} \;[+\;\text{geometría}],
\]
con variables estandarizadas; un coeficiente es el cambio en log-odds por una
desviación estándar (o cabeza contra pie) a igualdad del resto. Añadir la
geometría del remate mejora el AUC en \cifraDauc{} (IC por bootstrap por
partido con predicciones fijas). El JSON no trae el intercepto ni las escalas,
así que no se dibuja una curva en metros. Fuente: \fuente{balon\_parado\_v2}.}

\subsection*{3.4 · ¿Qué viaja con él y qué se queda en el club? \hfill{\normalfont\small diferenciador}}
\bloque{Una tabla-figura con cuatro eras (Jardine y Ortiz, en el América y fuera)
y la lista de técnicos con varios clubes.}
{Cada celda da el valor y una barra con su intervalo y la línea del cero. La
lectura es \nivel{C}: en el América los dos técnicos comparten posesión larga y
dominio territorial; fuera del América, el perfil de Jardine se invierte. Entre
los técnicos con varios clubes fuera del América, \cifraViajan{} conservan el
signo de su ajuste al marcador.}
{La comparación no separa plantel, presupuesto ni calendario, que cambian con
el club; con dos técnicos no hay regla general. Fuentes: \fuente{did\_h4\_v1},
\fuente{metricas\_v1}, \fuente{jugadores\_v1}, \fuente{contexto\_v1 ›
viaja\_marcador\_M1}.}

\subsection*{Cierre · ¿El método distingue de verdad? \hfill{\normalfont\small credibilidad}}
\bloque{Conteos crudos contra corregidos, el marcador de predicciones
preinscritas (\cifraMarcador) y el anexo de errores silenciosos.}
{Punto lleno: la predicción se cumplió; aro: falló. Tocar cada punto muestra su
enunciado. Los fallos se quedan a la vista.}
{Cada ADR escribió sus predicciones antes de correr el análisis. ADR-53 no dejó
su bloque de predicciones en el JSON; la página las recalcula con los cuatro
criterios de 06\_DECISIONS. El anexo de errores queda pendiente hasta tener su
fuente.}

\subsection*{Pendientes declarados}
2.2 (progresión) y, en el acto 1, dónde vive una posesión viva y la curva de
supervivencia llegan con ADR-61. 3.2 (plantel contra uso) y 3.3 (mapa de
estilos) llegan con ADR-60. La página los muestra como pendientes, sin cifras.

\subsection*{Cierre · Framework y límites \hfill{\normalfont\small 5.6}}
Lo que el marco no captura: la distribución de longitudes no es la de una
cadena de Markov; no hay tiempo en segundos, así que las transiciones van por
sustituto; no hay datos posicionales 360; hay competiciones fuera de los datos;
las fronteras de cada era tienen un margen de $\pm$1 partido. La simulación de
posesiones no se incluye. Nada es causal: el informe describe cómo jugaron los
equipos bajo cada técnico, contra la liga del mismo torneo; no mide
rendimiento ni intención.

%======================================================================
\section{Adenda 1 de ADR-59}\label{sec:adenda}
\begin{enumerate}
\item Los redondeos siguen al JSON: Jardine--Ortiz, extremo superior
  \cifraOrtJic; en contexto, minuto 60 con un decimal según el JSON.
\item Las predicciones de ADR-53 se recalculan desde \fuente{did\_h4\_v1}.
\item ADR-57 P1 se cuenta sin Jardine ni Ortiz (\cifraViajan); el JSON registra
  el conteo con los once técnicos, y la predicción falla con los dos.
\item Una frase B necesita intervalo. Pasan a \nivel{C}: la serie por torneo,
  la tabla de la liga en contexto, el embudo de balón parado y los cambios de
  la liga tras el primer cambio.
\item Figura 7: coeficientes estandarizados con IC en lugar de una curva en metros.
\item N80 ``alto'' = por encima de la mediana de la liga del torneo. Se fija
  después de ver los datos y así se declara.
\item El anexo de errores queda pendiente hasta tener el catálogo.
\item La página muestra $q$ con su lectura en palabras y nunca $p$.
\end{enumerate}

%======================================================================
\section{Adenda 2 de ADR-59}\label{sec:adenda2}
\begin{enumerate}
\item Tres actos y cinco historias en lugar de diez secciones fijas. El elenco
  se eligió después de ver el barrido exploratorio de eras; ninguna distancia
  del barrido entra a la página hasta ADR-60.
\item Las eras de una historia se buscan por igualdad exacta del nombre del
  técnico. La era principal es la de más partidos.
\item Las lecturas preinscritas de ADR-59 valen solo para Jardine y conservan
  sus guardas. Las otras historias usan plantillas fijas: \nivel{A} ``difiere''
  o nulo con margen, \nivel{B} ``por encima / por debajo de la liga'' o ``no se
  separa de la liga'', \nivel{C} conteos y rangos. Ninguna frase se elige por
  su resultado.
\item Presión: solo en los seis clubes de ADR-54. Una historia sin eras en ellos
  lo declara.
\item Cuatro capas por sección; la capa ``en un partido'' repite cifras de la
  capa 1 en unidades de un partido y no es una jugada real.
\item Portada con cinco ranuras fijas; interruptor sencilla / técnica.
\item Frases prohibidas nuevas: ``pesa más que'', ``demuestra que'', ``su idea''.
\end{enumerate}

\end{document}
